%==============================================================================
%  CAPÍTULO — TEORÍA GENERAL Y ANÁLISIS ECONÓMICO
%==============================================================================
\chapter{Teoría general y análisis económico del derecho concursal}
\label{cap:teoria-economica}

\fichaCapitulo{Por qué existe un derecho concursal y qué problema resuelve.}{El problema del
recurso común, la teoría del acuerdo hipotético entre acreedores, las tesis redistributivas, la
economía de la descarga y la justificación de las preferencias.}

%==============================================================================
\bloqueTeorico
%==============================================================================

\subsection{La pregunta fundacional: ¿por qué un derecho concursal?}

La existencia misma del derecho concursal no es evidente. Si el ordenamiento ya ofrece a cada
acreedor una acción de cobro y una ejecución individual, ¿por qué añadir un procedimiento
colectivo que suspende esas acciones, expropia la administración del deudor y redistribuye el
producto conforme a reglas propias?

La respuesta convencional ---\emph{para proteger al deudor}--- es insuficiente y, en rigor,
histórica más que dogmática: el concurso nació como instrumento de los acreedores, no del deudor.
La respuesta correcta exige identificar una \destacar{falla específica} que la ejecución individual
no puede resolver.

\subsection{El problema del recurso común (\emph{common pool problem})}

Esa falla es el \destacar{problema del recurso común}. Cuando el patrimonio del deudor es
insuficiente para satisfacer a todos, cada acreedor tiene un incentivo individual a ejecutar
primero, porque quien llega antes cobra íntegro y quien llega después no cobra nada. El resultado
agregado de esas decisiones individualmente racionales es colectivamente destructivo:

\begin{enumerate}
  \item \textbf{Carrera de ejecuciones.} Los acreedores gastan recursos en llegar primero
  ---costos de litigación que no crean valor, sino que sólo determinan quién se lo lleva---.

  \item \textbf{Desmembramiento del activo.} La ejecución individual realiza bienes
  \emph{aisladamente}. Si el conjunto vale más que la suma de sus partes ---una fábrica en marcha
  frente a sus máquinas subastadas por separado---, la carrera destruye ese excedente de valor de
  empresa en funcionamiento.

  \item \textbf{Liquidación prematura.} Empresas económicamente viables, con problemas meramente de
  liquidez, son desmanteladas porque ningún acreedor individual tiene incentivo a esperar mientras
  los demás ejecutan.
\end{enumerate}

\begin{foco}[La analogía clásica del recurso común]
El problema es idéntico al de un lago sobreexplotado o un yacimiento con múltiples concesionarios:
cada pescador sabe que la captura es insostenible, pero si él se abstiene y los demás no, sólo
pierde. La solución no es moral sino institucional: una regla que suspenda la extracción individual
y organice la explotación colectiva. El concurso \emph{es} esa regla aplicada al patrimonio del
deudor insolvente, y de ahí derivan sus dos efectos característicos: la
\destacar{paralización de las ejecuciones individuales} (\art{57} y \art{135}) y la
\destacar{administración unificada} de la masa.
\end{foco}

\subsection{La teoría del acuerdo hipotético entre acreedores}

La formulación más influyente del punto anterior es la del \emph{creditors' bargain}: el derecho
concursal debe entenderse como el \destacar{contrato que los acreedores habrían celebrado} de haber
podido negociar entre sí, sin costos de transacción, antes de que sobreviniera la insolvencia.

Bajo ese supuesto, ningún acreedor racional habría aceptado una carrera destructiva de valor: todos
habrían convenido de antemano en suspender las ejecuciones individuales, someterse a una
administración común y repartirse el producto conforme a las prelaciones pactadas
\emph{ex ante}. El concurso no crea derechos nuevos: se limita a hacer cumplir colectivamente lo
que las partes habrían acordado.

De esta tesis derivan dos corolarios de enorme importancia práctica:

\paragraph{Corolario de neutralidad.} El concurso debe \destacar{respetar las prelaciones
sustantivas} preexistentes. Si un acreedor negoció una hipoteca, su rango debe conservarse dentro
del procedimiento; alterarlo equivaldría a expropiar en el concurso lo que se adquirió fuera de él.
Es el fundamento teórico de la conservación de preferencias del \art{61} y del respeto al orden del
\artcc{2472} que impone el \art{241}.

\paragraph{Corolario de maximización.} La finalidad del procedimiento es
\destacar{maximizar el valor de la masa}, y la elección entre reorganizar y liquidar es una
cuestión económica, no moral: se reorganiza si la empresa vale más en marcha que desmantelada, y se
liquida en caso contrario. La conservación de la empresa no es un fin en sí mismo.

\subsection{Las tesis redistributivas y la crítica al modelo contractual}

La teoría del acuerdo hipotético ha sido objeto de una crítica sustancial que conviene exponer con
rigor, porque explica buena parte del diseño de la \LIR{}.

La objeción central es que el modelo contractual \destacar{sólo considera a los acreedores
negociadores} y deja fuera a quienes no pudieron negociar nada: los trabajadores, los acreedores
extracontractuales ---la víctima de un daño no eligió a su deudor---, los proveedores pequeños sin
poder de mercado y la comunidad afectada por el cierre de una empresa. Si el concurso fuera
únicamente el contrato hipotético de los acreedores financieros, esos intereses carecerían de toda
tutela.

Desde esta perspectiva, el derecho concursal es un \destacar{mecanismo de distribución de las
pérdidas} de un fracaso económico, y esa distribución es una decisión política que la ley adopta y
que no se deduce de ningún acuerdo hipotético.

\begin{alerta}[La \LIR{} es un compromiso entre ambas teorías]
El derecho chileno no adopta ninguna de las dos posiciones en estado puro, y advertirlo evita
lecturas dogmáticas erróneas. Es \emph{contractualista} cuando conserva las preferencias
(\art{61}), respeta la prelación civil (\art{241}) y somete la continuación del giro a un criterio
de mayor recuperación. Es \emph{redistributivo} cuando otorga superpreferencia a los créditos
laborales (\artcc{2472} \Nro 5 y 8), cuando sustrae a los trabajadores del acuerdo de
reorganización (\art{60 A}) y cuando concede al deudor persona natural una descarga que ningún
acreedor habría convenido. Identificar en qué registro opera cada norma es la clave para
interpretarla correctamente.
\end{alerta}

%==============================================================================
\bloqueNormativo
%==============================================================================

\subsection{La economía de la descarga}

La extinción de saldos insolutos es, desde la lógica del acuerdo hipotético, una anomalía: ningún
acreedor consentiría en renunciar a su crédito. Su justificación debe buscarse fuera de ese marco,
y admite tres lecturas complementarias que el capítulo \ref{cap:evolucion} ya anticipó y que aquí
se sistematizan.

\paragraph{La descarga como seguro.} El deudor no puede diversificar el riesgo de su propio fracaso
económico; el acreedor profesional, en cambio, presta a miles de deudores y \emph{sí} puede
diversificarlo y tarificarlo en la tasa de interés. La descarga asigna el riesgo residual a quien
lo administra a menor costo, funcionando como un seguro social obligatorio financiado por el
sistema crediticio.

\paragraph{La descarga como incentivo.} La deuda perpetua e impagable opera como un
\destacar{impuesto marginal del cien por ciento} sobre el ingreso futuro del deudor: si todo lo que
gane será absorbido por acreedores preexistentes, carece de motivo racional para generar ingreso
formal. El resultado es el desplazamiento hacia la informalidad, que perjudica al deudor, al fisco
y ---paradójicamente--- a los propios acreedores, que no recuperan de un patrimonio que no se
genera. La descarga restituye el incentivo a producir.

\paragraph{La descarga como protección del capital humano.} A diferencia de la sociedad, la persona
natural no puede disolverse. Su principal activo no es patrimonial sino su
\destacar{capacidad de generar ingresos}, y un endeudamiento perpetuo la inutiliza. La descarga
preserva ese activo en beneficio del conjunto.

\begin{foco}[Por qué la descarga tiene un filtro de probidad y no de generosidad]
Las tres justificaciones anteriores presuponen un fracaso \emph{genuino}. Si el deudor provocó su
insolvencia o la simuló, ninguna opera: no hay riesgo que asegurar, sino un riesgo creado; no hay
incentivo que restituir, sino un premio al oportunismo. De ahí que todo sistema de descarga incluya
un filtro ---el \art{169 A} chileno, la buena fe española, el \S~523 estadounidense--- y que ese
filtro sea la pieza que hace defendible el conjunto. La descarga no es indulgencia: es una
herramienta económica que sólo funciona bien calibrada.
\end{foco}

\subsection{La justificación económica de las preferencias}

El orden de prelación del \artcc{2472} no es arbitrario, y cada categoría responde a una lógica
distinta que conviene explicitar:

\begin{description}[leftmargin=0pt,style=nextline,font=\sffamily\bfseries]
  \item[Gastos del procedimiento (N.\textsuperscript{o} 1 y 4)] Sin prioridad para quien financia o
  administra el concurso, nadie prestaría el servicio ni aportaría recursos: la masa no se
  administraría. Es la misma lógica del \emph{new money} del \art{74}. La preferencia es aquí
  puramente \emph{instrumental}.

  \item[Créditos laborales (N.\textsuperscript{o} 5 y 8)] El trabajador es un
  \destacar{acreedor involuntario y no diversificado}: no eligió financiar a la empresa, no pudo
  tarificar el riesgo, concentra en ese único deudor la totalidad de su ingreso y carece de
  información sobre la solvencia. Todas las razones que justifican asignar el riesgo al acreedor
  profesional operan aquí en sentido inverso.

  \item[Créditos del fisco por retención y recargo (N.\textsuperscript{o} 9)] Se trata de sumas que
  el deudor \destacar{recaudó de terceros} ---el IVA que cobró, las retenciones que descontó--- y no
  enteró. No son dinero suyo: la preferencia corrige una apropiación, no privilegia al fisco.

  \item[Garantías reales (segunda y tercera clase)] Su respeto no es un privilegio, sino la
  condición de existencia del crédito garantizado: si la garantía pudiera desconocerse en el
  concurso ---precisamente cuando importa---, dejaría de cumplir función alguna y el costo del
  crédito subiría para todos los deudores.
\end{description}

\begin{practica}[Usar la justificación al litigar la preferencia]
Cuando se discute el alcance de una preferencia ---si tal partida laboral entra en el
N.\textsuperscript{o} 5 o en el 8, si un tributo es de retención o de declaración---, el argumento
teleológico es decisivo y suele estar ausente de los escritos. Mostrar que la interpretación
propuesta \emph{sirve a la razón de ser} de la categoría ---proteger al acreedor no diversificado,
o restituir lo recaudado de terceros--- es más persuasivo que el mero cotejo literal, y los
tribunales lo reciben bien.
\end{practica}

\subsection{Reorganizar o liquidar: el criterio económico}

La decisión entre reorganización y liquidación admite una formulación precisa. Debe reorganizarse
cuando el \destacar{valor de la empresa en marcha} supera su \destacar{valor de liquidación}; debe
liquidarse en el caso inverso. La diferencia entre ambos ---el valor de empresa en funcionamiento---
es el excedente que el procedimiento debe preservar y repartir.

De ahí se siguen dos precisiones que la práctica confunde con frecuencia:

\paragraph{Insolvencia no es inviabilidad.} Una empresa puede ser insolvente ---pasivo superior al
activo--- y económicamente viable: su operación genera márgenes positivos, pero arrastra una
estructura de deuda desproporcionada. El instrumento correcto es la reorganización, que ajusta el
\emph{pasivo} sin tocar la operación. A la inversa, una empresa puede ser solvente pero inviable, y
su liquidación ordenada es la solución eficiente.

\paragraph{El sesgo de continuación.} Los administradores y los socios tienen un incentivo
sistemático a prolongar una empresa inviable, porque juegan con dinero ajeno: si la apuesta
resulta, recuperan su participación; si fracasa, la pérdida adicional la soportan los acreedores.
Esta asimetría ---apostar por la resurrección--- justifica los controles de la reorganización: la
supervisión del veedor, los límites de disposición del \art{74} y las mayorías cualificadas del
\art{79}.

\begin{alerta}[Los números chilenos y la hipótesis del acceso]
Las cifras del sistema chileno ---más del 93\,\% de procedimientos de persona deudora y apenas
unas decenas de reorganizaciones empresariales al año (capítulo \ref{cap:renegociacion})--- admiten
dos lecturas opuestas. La optimista: casi no hay empresas viables en crisis. La realista: el
instrumento de reorganización es tan costoso y su resultado tan incierto ---sin arrastre de clases,
el rechazo de una minoría lo frustra--- que las empresas viables sencillamente no lo usan y
negocian fuera del sistema o mueren sin procedimiento. La comparación con los ordenamientos que sí
disponen de arrastre entre clases (capítulos \ref{cap:comparado-eeuu} y \ref{cap:comparado-espana})
inclina la balanza hacia la segunda lectura.
\end{alerta}

\subsection{Los costos del concurso}

Todo procedimiento concursal genera costos que reducen lo que finalmente reciben los acreedores, y
distinguirlos ayuda a evaluar cualquier propuesta de reforma:

\begin{description}[leftmargin=0pt,style=nextline,font=\sffamily\bfseries]
  \item[Costos directos] Honorarios del liquidador o veedor, costas, tasaciones, publicaciones,
  auditorías. Son visibles y mensurables, y la tabla del \art{40} los administra explícitamente.

  \item[Costos indirectos] Pérdida de clientes y proveedores ante el anuncio del concurso, fuga de
  personal clave, deterioro de activos durante la tramitación, pérdida de oportunidades de negocio.
  Suelen ser \destacar{muy superiores} a los directos y no aparecen en ninguna cuenta.

  \item[Costos \emph{ex ante}] El diseño concursal afecta el costo del crédito \emph{antes} de
  cualquier insolvencia: un régimen que debilite las garantías encarece el crédito garantizado para
  todos los deudores, incluidos los que nunca caerán en concurso.
\end{description}

\begin{foco}[Por qué la celeridad es una cuestión sustantiva y no de mera eficiencia]
Los costos indirectos crecen con el tiempo: cada mes de tramitación erosiona el valor de la masa.
Esto convierte la celeridad ---la simplificación procedimental, la restricción recursiva del
\art{4}, los plazos breves--- en una decisión de \emph{fondo} sobre la protección del patrimonio, y
no en una preferencia administrativa. Es también el argumento que el Tribunal Constitucional acogió
al validar la limitación de recursos (capítulo \ref{cap:liquidacion-ordinaria}): la celeridad
protege a la masa, y por eso puede justificar restricciones procesales que en otro contexto serían
sospechosas.
\end{foco}
